% Draft v1 of §4.2 (lines 40-54 of sections/eval-on-error-distance.tex), 2026-11-02.
% Numbers: pro-paper.seen artifacts, plus the S-P taxonomy prototype
% (tasks/20261102-sp-taxonomy-module) for the deviation paragraph. Those are
% marked \rednote{taxonomy} until report-sp-taxonomy exists. All numbers are
% from before the min-VPs-per-TG filter and shift slightly once it lands.

\subsection{Understand the Baseline: Shortest Ping}
\label{sec:err-characterization-sp}

S-P is the most transparent of the six methods. It returns the location of the VP with the lowest RTT to the TG, so its answer is always a VP location. Our VP set is dense, yet S-P classifies only half of the TGs correctly (\S~\ref{sec:eval-error-distance}). This subsection identifies what limits it.

\noindent \textbf{How a measurement reaches a TG.} Fig.~\ref{fig:network-topology} illustrates the path of a measurement. Traffic from a VP leaves the MNO, enters the TG's network at an interconnect $X$, and travels inside that network to the TG. Its propagation distance is approximately $D_{prop} = D_{VP\text{-}X} + D_{X\text{-}TG}$, the sum of the great-circle distances from the VP to $X$ and from $X$ to the TG. The BGP announcements of the TG's network decide which interconnect the traffic enters through, and they change infrequently~\citeneeded{stability of BGP ingress / interconnect selection}. We call this interconnect the TG's \emph{ingress point} $X_{\mathrm{in}}$. It is usually the interconnect nearest the TG, since content providers optimize for latency~\citeneeded{latency-driven ingress selection}. VPs colocated with $X_{\mathrm{in}}$ have $D_{VP\text{-}X} \approx 0$, so they measure the lowest RTTs to the TG and form its \emph{low-RTT tier}. S-P returns a VP of this tier, so it answers near $X_{\mathrm{in}}$ instead of near the TG, and its error distance is approximately $D_{X\text{-}TG}$.

Fig.~\ref{fig:sp-scatter} tests this. It plots S-P's error against the TG's distance to its nearest VP (Fig.~\ref{fig:sp-vp-scatter}) and to its nearest interconnect, which we take as $X_{\mathrm{in}}$\footnote{We did not record the ingress point during measurement, so we use the TG's nearest interconnect in its place. On every TG, the TG's own cell holds its nearest interconnect exactly when it holds any interconnect.} (Fig.~\ref{fig:sp-x-scatter}). S-P's error follows the distance to the interconnect (Spearman's $\rho = 0.86$) much more closely than the distance to the nearest VP ($\rho = 0.38$). S-P's answer also lands in the cell of the TG's nearest interconnect for 85\% of the TGs, against 6.7\% for a VP drawn at random from those that measured the TG. We explain the points off the diagonal below. \rednote{calibrate data}

\noindent \textbf{Whether the ingress point lies in the TG's cell explains most of S-P's accuracy.} The answer space (\S~\ref{sec:answer-space}) lets us test this directly. We call a TG \emph{has-$X$} if its own cell holds an interconnect, and \emph{no-$X$} otherwise. Since S-P answers near $X_{\mathrm{in}}$, it can classify a TG correctly only when $X_{\mathrm{in}}$ lies in the TG's own cell. Pooled over the three networks, 53\% of the TGs are has-$X$ and 47\% are no-$X$. S-P classifies 93\% of the has-$X$ TGs correctly and almost none of the no-$X$ TGs ($<$1\%; Fig.~\ref{fig:sp-accuracy-by-x}). Its wrong no-$X$ answers land in the cell next to the TG's in 88\% of the cases. Per network, S-P's accuracy matches the has-$X$ share in AS-B (42\% vs.\ 42\%) and AS-C (43\% vs.\ 43\%), and falls 11~pp short of it in AS-A (64\% vs.\ 75\%), where all of its has-$X$ misses occur (Table~\ref{tab:x-cell}). The no-$X$ TGs do not fail for lack of nearby VPs: their nearest VP is as close as that of the has-$X$ TGs (median 3.6 vs.\ 3.0 $\times 10^{-3}$). They fail because their low-RTT tier sits at an $X_{\mathrm{in}}$ outside their cell. Their minimum RTTs are accordingly higher (Fig.~\ref{fig:boxplot-rtt-of-x-cell-by-asn}), with medians of 0.095--0.148 across the three networks, 5.4--9.6$\times$ those of the has-$X$ TGs. Every no-$X$ TG lies above 0.032 (the dotted line), while 92\% of the has-$X$ TGs lie below it. \rednote{calibrate data}

\noindent \textbf{Non-discriminative RTTs explain the has-$X$ misses.} S-P relies on the minimum RTT to single out the low-RTT tier, and within it the VP nearest $X_{\mathrm{in}}$. When candidate VPs measure RTTs too close to tell apart, S-P can pick a VP outside the TG's cell. Two cases account for all of the 7\% of has-$X$ TGs that S-P misses, both in AS-A. \rednote{taxonomy}

(1) \emph{Near-ties within the tier} (3.0\% of has-$X$ TGs). With the speed of the Internet, VPs within 100~km of the same interconnect differ by at most about 1~ms in propagation RTT, which a small difference in access delay can outweigh. For these TGs, the VPs within 1~ms of the minimum RTT span the TG's own cell and a neighbouring one, and S-P picks the neighbour. This never happens on a has-$X$ TG that S-P classifies correctly: there, every VP within 1~ms of the minimum lies in a single cell. The same effect spreads the lower left of Fig.~\ref{fig:sp-x-scatter}, but there it rarely moves the answer out of the TG's cell. \rednote{taxonomy}

(2) \emph{Inflated RTTs} (3.7\% of has-$X$ TGs). Every VP measures a heavily inflated minimum RTT to these TGs, 0.64--0.69 against at most 0.10 for any correctly classified has-$X$ TG, because of a deliberate routing policy we confirmed internally. A delay shared by every VP then swamps the distance signal: the RTT no longer correlates with VP distance (median $\rho \approx 0$), and VPs in several cells measure within 1~ms of the minimum. S-P's answers fall far from the TGs; these are the outliers near $10^3$ in Fig.~\ref{fig:sp-x-scatter}. The same policy affects 3.3\% of the no-$X$ TGs, which S-P misses either way. \rednote{taxonomy}

\noindent \textbf{When the ingress point is not the nearest interconnect.} In 17\% of the no-$X$ TGs, all in AS-B and AS-C, the RTTs are inconsistent with traffic entering at the TG's nearest interconnect. For 9.3\% of the no-$X$ TGs, some VP measures an RTT below the propagation floor of any path through that interconnect, while another interconnect fits. For a further 7.8\%, another interconnect explains the RTTs better. In both cases S-P answers at that other interconnect, which is consistent with S-P locating the actual ingress point. This changes where S-P's answer lands but not whether it is correct, since a no-$X$ TG's cell holds no interconnect to land in. \rednote{taxonomy}

\noindent \textbf{Takeaways.} \textit{In an operator's setting, at least one VP sits at every interconnect of the TG's network, so the low-RTT tier of VPs sits at the TG's ingress point and S-P answers there. S-P classifies a TG correctly when that ingress point lies in the TG's serving region. It fails when RTTs stop discriminating between VPs, because they are nearly tied or inflated by a delay every VP shares.}
